\documentclass[a4paper,dvipsnames, 11pt]{article}
\usepackage{latexsym}
\usepackage{amsmath}
%\usepackage{mathspec}
%\usepackage{amsthm}
% \usepackage{amssymb}
% \usepackage{textcomp} 
% \usepackage{ulem}
% % \usepackage{fontspec}     
\usepackage[british]{babel}
\usepackage{enumitem}
% \usepackage{pifont}
\usepackage[a4paper, margin=1.2in]{geometry}
\usepackage{hyperref}
\usepackage[dvipsnames, table]{xcolor}
\usepackage{tikz}
% \usepackage{mathrsfs}

% \usepackage{eucal}
\usepackage{dsfont}
% \usepackage[most]{tcolorbox}
% \usepackage{mathdots}
\usepackage{minted} %? used for code-integration
% \usepackage{titling}
\usepackage[bibitemsep=100pt, citestyle=alphabetic,bibstyle=alphabetic]{biblatex}

\definecolor{darkmagenta}{RGB}{99,0,99}
\hypersetup{
    colorlinks,
    citecolor=darkmagenta,
    filecolor=darkmagenta,
    linkcolor=darkmagenta,
    urlcolor=darkmagenta
} %? Change hyperlink's styles

%%%%%%%%%%%%%%? courbes
\usepackage{amsfonts,amscd}
\usepackage{pgfplots}
%%%%%%%%%%%%%?

%%%%%%%%%%%%%
\DeclareMathOperator{\mat}{Mat}
\DeclareMathOperator{\card}{Card}
\DeclareMathOperator{\len}{len}
\DeclareMathOperator{\erf}{erf}
%%%%%%%%%%%%%


%%%%%%%%%%%%%
\newcommand{\dt}{\mathrm{d}}
\newcommand{\un}{\mathds1}
\renewcommand{\o}{\scriptstyle\mathcal{O}}
\renewcommand{\O}{\mathcal{O}}
\renewcommand{\emptyset}{\varnothing}
\newcommand{\lbd}{\lambda}
\renewcommand{\phi}{\varphi}
\def\inte #1 #2 { [\![#1,#2]\!] }
%%%%%%%%%%%%%


%%%%%%%%%%%%%
\newcommand{\N}{\mathbb{N}}
\newcommand{\R}{\mathbb{R}}
\newcommand{\Z}{\mathbb{Z}}
\newcommand{\C}{\mathbb{C}}
\newcommand{\K}{\mathbb{K}}
\newcommand{\Q}{\mathbb{Q}}
\renewcommand{\P}{\mathbb{P}}
\newcommand{\M}{\mathcal{M}}
\renewcommand{\r}{\mathcal{R}}
\renewcommand{\S}{\mathcal{S}}
%%%%%%%%%%%%%


\renewcommand{\arraystretch}{1.4} %? For better matrix
\renewcommand{\thesection}{\arabic{section}} %? Roman number for sections


\newenvironment{code}{
    \VerbatimEnvironment
    \begin{minted}[mathescape,
        linenos,
        numbersep=5pt,
        gobble=2,
        frame=lines,
        framesep=2mm]{cpp}%
}{
        \end{minted}%
}

    \setcounter{secnumdepth}{4}
    \setcounter{tocdepth}{4}

\hbadness=10000
\vbadness=10000
\flushbottom
\addbibresource{rapport.bib}
\usepackage{tkz-base}
\usepackage{algorithm}
\usepackage{algorithmic}
\usepackage{wrapfig}
\usepackage{subcaption}
\usepackage{changepage}
\usepackage[]{tocbibind}
\usepackage[useregional]{datetime2}
\usepackage{adjustbox}
\usepackage{pdfpages}
\usepackage{tabularx}
\usepackage{array}
\usepackage{cleveref}
\usepackage{multicol}
\usepackage{graphicx}
\usepackage{multirow}
\usepackage[multiple]{footmisc}
\usepackage[minNoteWidth=1cm,
    colorinlistoftodos,
    backgroundcolor=magenta!30!white,
    bordercolor=magenta]{luatodonotes}
\usepackage{amssymb}

\newcolumntype{C}{>{\centering\arraybackslash}X}
\newcolumntype{s}{>{\hsize=.3\hsize\linewidth=\hsize}C}
\newcolumntype{D}{>{\hsize=.4\hsize\linewidth=\hsize}C} %Double width column

\setlength{\bibhang}{0pt}

\defbibenvironment{bibliography}
{\list
{\printtext[labelalphawidth]{\printfield{labelprefix}\printfield{labelalpha}}}% Updated line
{\setlength{\labelwidth}{\labelalphawidth}% Updated line
\setlength{\leftmargin}{0pt}%
\setlength{\labelsep}{\biblabelsep}%
\setlength{\itemsep}{\bibitemsep}%
\setlength{\itemsep}{0pt}% Set to 0pt to remove vertical space
\setlength{\parsep}{\bibparsep}}%
\renewcommand*{\makelabel}[1]{\hss\hspace{\dimexpr\labelalphawidth+\labelsep}##1}}
{\endlist}
{\item}

% Redefine the footnotetext command to remove indentation
\makeatletter
\renewcommand\@makefntext[1]{%
    \noindent\makebox[0em][r]{\@thefnmark\hspace{5pt}}#1}
\makeatother

\newcommand{\reftodo}[1][]{\todo[backgroundcolor=green!30!white, bordercolor=green, #1]{\thesubsubsection{}~Ref} {\color{red} REF}}
\newcommand{\missingtext}[1][]{\todo[inline, caption=\thesubsection~Fill missing text, backgroundcolor=cyan!30!white, bordercolor=cyan]{#1}}
%\usepackage[]{todonotes}

\setlength\parindent{0pt}

\usetikzlibrary {shapes.multipart}

\makeatletter
\newenvironment{fixedfigure}
{\def\@captype{figure}\center}
{\endcenter}
\makeatother



% \usepackage{graphicx,txfonts}

\addto\captionsbritish{
    \renewcommand{\contentsname}%
    {Table of content}%
}

\renewcommand{\today}{\number\day\textsuperscript{\scriptsize\text{th}} \DTMenglishmonthname{\month} \number\year}


\makeatletter
\newcommand\setxveclength[5]{% newmacro, node1, node2
    \pgfpointdiff{#2}{#3}
    \edef#1{\the\pgf@x}
}
\makeatother

\sloppy                  

\pgfplotsset{compat=1.16}

\begin{document}


%\setmathfont{Latin Modern Math}
%\setmathfont[range={\mathscr,\mathbfscr}]{XITS Math}

% \maketitle



\begin{titlepage}
    \begin{center}
        \vspace*{1cm}
            
        \Huge
        \textbf{Random Osborne Algorithm for Matrix Balancing}
            
        \vspace{0.5cm}
        \LARGE
        Optimal transport report
            
        \vspace{2.5cm}
            
        \huge {\sc Alexi Canesse\footnote{\href{mailto:alexi.canesse@ens-lyon.fr}{alexi.canesse@ens-lyon.fr}, ENS de Lyon, France}\footnote{\href{mailto:alexi.canesse@ens-paris-saclay.fr}{alexi.canesse@ens-paris-saclay.fr}, ENS Paris-Saclay, France}}\\\Large
            
        \vfill

        Optimal transport course\\
        Part of the MVA program at ENS Paris-Saclay.
            
        \vspace{2cm}
        \includegraphics[width=0.5\textwidth]{"./figures/logo.pdf"}
        \vspace{2cm}

        \Large
        Computer Science and Mathematics\\
        École Normale Supérieur Paris-Saclay\\
        Orsay, France\\
        \today
            
    \end{center}
\end{titlepage}





\newpage

%\tableofcontents
\makeatletter
\section*{Table of content}
\vspace{.05\textwidth}
\@starttoc{toc}
\thispagestyle{empty} % Remove page number from the ToC page
\clearpage % Start the new page without page number

\setcounter{page}{1} % Reset page counter and set page number to 1

\makeatother
\newpage

\vspace*{10pt}
\begin{adjustwidth}{30pt}{30pt}
    \textit{\todo{Abstract (½ page): 
    What problem(s) is studied? 
    Why is it relevant? 
    What solution(s) is proposed? 
    Which contributions (theory, numerics, etc) ? }
        Accurate physics simulations can be computationally expensive due to the requirement of fine meshes.
        However, a cost-effective approach involves reducing the mesh size while retaining crucial features required
        for the simulation.
        Some of these features are conveniently represented in the spectra of Laplacians operators defined on the mesh.
        There exists a coarsening method that filters out unnecessary spectral bands.
        The challenge lies in determining which spectral bands are important for the simulation, as \textbf{this information is
        largely unknown}.
    }\\

    \textit{
        To address this, we propose a \textbf{novel approach} to guide the coarsening process using \textbf{insights
        from physical simulations}.
        By learning the spectral subspaces that significantly influence the simulation outcomes, we can create a coarse
        mesh that retains the characteristics essential for an accurate simulation.
        Main contributions are:
        \begin{itemize}
            \item A novel spectral optimization method for coarsening, with physical simulations in mind.
            \item Understanding importance of spectral subspaces: By understanding the spectral subspaces that contribute
            the most to the simulations, we gain valuable insights into the underlying physics and behavior of the system.
            \item Faster simulation times: The reduced mesh size leads to a considerable speed-up in the simulation
            process, making it more efficient and feasible for practical applications.
        \end{itemize}
    }
\end{adjustwidth}
\vspace*{10pt}

\section{Introduction (3 pages)}

\subsection{Presentation of the problem}

\subsection{Related work}
Previous works (at least a few citations). If relevant, include things that you have seen during the MVA course (or possibly other courses). 

\subsection{Contributions}
Contributions. Why is the studied method different/better/worse/etc. than existing previous works. 

\section{Main body (10 pages)}

\subsection{Presentation of the method}
Include variants etc

\subsection{Theoretical guarantees}

\subsection{Numerics}

\section{Conclusion and perspective}

Summary of the result obtained: pros and cons (limitation, problems, error in the articles, etc)
Possible improvement/extension

\section{Connexion with the course}

MANDATORY SECTION:. What are the notions/results/algorithms presented in the course that are used or related to the one presented in this paper?

\end{document}


